% Braccio teso in orizzontale con il manubrio (problema ud2-rigido-deltoide):
% la spalla O fa da centro dei momenti. Il deltoide tira a 15 cm dalla spalla,
% verso l'alto e verso il corpo, a 30 gradi dal braccio; la sua retta d'azione
% (tratteggiata) passa a b_d = 15 cm sin 30 = 7,5 cm dalla spalla. Il peso del
% braccio agisce a 30 cm, quello del manubrio a 60 cm. Le distanze sono in
% scala, le frecce non sono in scala con i moduli delle forze.
\begin{tikzpicture}[>=Stealth, font=\small, x=0.15cm, y=0.15cm]
  \coordinate (O) at (0,0);
  \coordinate (D) at (15,0);
  \coordinate (H) at (3.75,6.5);
  % braccio
  \filldraw[fill=esempio!20, draw=esempio, thick, rounded corners=3pt] (-1,-1.2) rectangle (61,1.2);
  \fill[richiamo!70] (O) circle (1.1);
  \fill (O) circle (1.5pt);
  \node[left, align=right] at (-1.2,0) {spalla\\$O$};
  % deltoide e sua retta d'azione
  \draw[dashed, richiamo] (D) -- ($(D)!1.4!(H)$);
  \draw[->, medica!80!black, very thick] (D) -- ++(150:10) node[above right] {$\vec F_d$};
  \draw ($(D)+(180:6)$) arc (180:150:6);
  \node at ($(D)+(165:8.6)$) {\ang{30}};
  % braccio del deltoide, perpendicolare alla retta d'azione
  \draw[formula, thick] (O) -- (H) node[midway, left] {$b_d$};
  \draw[formula] ($(H)!0.15!(O)$) -- ++(-30:1.1) -- ($(H)!0.085!(D)$);
  % pesi
  \draw[->, attenzione, very thick] (30,-1.2) -- (30,-8) node[right] {$\vec F_b$};
  \filldraw[fill=richiamo!60, draw=black] (58.2,-1.2) rectangle (61.8,-4);
  \draw[->, attenzione, very thick] (60,-4) -- (60,-9) node[left] {$\vec F_m$};
  % distanze dalla spalla
  \draw[dashed, richiamo] (O) -- (0,-16);
  \draw[|<->|, formula] (0,-4.5) -- (15,-4.5) node[midway, below] {\SI{15}{\centi\metre}};
  \draw[|<->|, formula] (0,-10.5) -- (30,-10.5) node[midway, below] {\SI{30}{\centi\metre}};
  \draw[|<->|, formula] (0,-15) -- (60,-15) node[midway, below] {\SI{60}{\centi\metre}};
\end{tikzpicture}
